Este trabalho aborda o treinamento por aprendizado por reforço em larga escala no futebol de robôs IEEE Very Small Size Soccer (VSSS) 3v3 e faz três contribuições. Primeiro, um simulador em JAX que executa milhares de partidas em paralelo em um único acelerador. Sua física de contato modela efeito (\emph{spin}), deslizamento e saltos da bola nas paredes, detecta o primeiro contato da bola com paredes, traves, robôs e linha de gol por testes varridos e projeta iterativamente contatos acoplados; a lógica de árbitro e reinícios é compartilhada por três interfaces de ambiente e protegida por uma suíte de regressão de falhas de colisão. Segundo, uma suíte de treinamento Proximal Policy Optimization (PPO) totalmente no acelerador, com observação, recompensa e política compartilhada no estilo de Brand\~ao treinada com aumento de dados por simetria espelhada, autojogo histórico contra um oponente congelado e uma geração compilada que contém coleta, estimação de vantagem e otimização. Terceiro, um currículo \emph{online} baseado em retornos, que reapresenta situações iniciais conforme os retornos dos episódios recentes dentro do laço compilado. O simulador atinge \textbf{[X]} passos de ambiente por segundo em \textbf{[GPU]}, contra 63--378 passos por segundo por agente de treinamento nas configurações anteriores do ITAndroids baseadas em ODE; o ganho vem tanto da física 2D mais simples quanto da GPU, e os valores em CPU e GPU são reportados separadamente. Em uma ablação com orçamentos iguais, o currículo \textbf{[melhorou / não melhorou]} as políticas finais em relação a reinícios aleatórios uniformes (\textbf{[resultado]}). A principal evidência é qualitativa: comportamentos cada vez mais complexos (chute, bloqueio, posicionamento para segundas bolas, diferenciação de papéis) surgem ao longo das gerações. A equipe também é avaliada quantitativamente contra oponentes parados, aleatórios e versões anteriores de si mesma e, para a última política de cada linhagem, contra a estratégia clássica do ITAndroids.
